%=============================================================================
% APPENDIX W: EXPERIMENTAL PROTOCOLS AND SENSITIVITY ANALYSES
% Comprehensive pre-registered protocols for DFD tests
%=============================================================================

\section{Experimental Protocols and Sensitivity Analyses}\label{app:protocols}

This appendix provides detailed, pre-registered experimental protocols for 
the key DFD discriminators. Each protocol specifies the observable, 
prediction, systematics budget, decision rule, and falsification criteria.

%-----------------------------------------------------------------------------
\subsection{Cavity-Atom LPI Test: Complete Protocol}\label{app:proto-cavity}
%-----------------------------------------------------------------------------

The height-separated cavity--atom comparison remains a valuable protocol, but after geometric cancellation (Sec.~\ref{sec:cavity-proof}) it is best viewed as a demanding long-horizon residual test.  The tree-level cavity/atom response cancels; only a screened residual $\xi_{\rm LPI}^{\rm res}$ survives.  This section preserves the full protocol details for completeness and for future experiments that may reach the required sensitivity.

\subsubsection{Observable and Predictions}

The frequency ratio at height $h$ is:
\begin{equation}
R(h) \equiv \frac{\nu_C(h)}{\nu_A(h)}.
\end{equation}

\paragraph{GR prediction.}
\begin{equation}
\left(\frac{\Delta R}{R}\right)_{\mathrm{GR}} = 0.
\end{equation}

\paragraph{Corrected DFD prediction.}
After the constitutive-chain cancellation, the surviving signal is:
\begin{equation}
\left(\frac{\Delta R}{R}\right)_{\mathrm{DFD}} = \xi_{\rm LPI}^{\rm res}\,\frac{g\,\Delta h}{c^2},
\end{equation}
where $\xi_{\rm LPI}^{\rm res}$ is the screened residual coupling that remains once the leading geometric (tree-level) effect is removed.  At Earth's surface, the screening analysis of Sec.~\ref{sec:clocks-regime} and the BACON constraints of Sec.~\ref{sec:cavity-bacon} restrict $\xi_{\rm LPI}^{\rm res}$ to be small --- far below the order-unity value assumed in earlier internal drafts.

\paragraph{Numerical estimate.}
For $\Delta h = 100$ m and $g = 9.8$ m/s$^2$:
\begin{equation}
\frac{g\,\Delta h}{c^2} \simeq 1.1 \times 10^{-14}.
\end{equation}
The DFD signal is this factor \emph{multiplied by} the small screened residual $\xi_{\rm LPI}^{\rm res}$, making the target signal extremely demanding.  This is why the cavity--atom channel is ranked below cross-species and nuclear-clock tests in the current experimental priority ordering (Sec.~\ref{sec:clocks-priorities}).

\subsubsection{Experimental Configuration}

\begin{itemize}
\item \textbf{Lower station} at $h_1$: High-stability optical cavity (ULE or Si) 
and reference atomic clock (Sr or Yb lattice clock).
\item \textbf{Upper station} at $h_2 = h_1 + \Delta h$: Second atomic clock 
and auxiliary diagnostics.
\item \textbf{Link:} Phase-stabilized optical fiber at $< 10^{-18}$ level.
\item \textbf{Height difference:} $\Delta h \sim 100$ m (tower, elevator shaft, or mine).
\end{itemize}

\subsubsection{Measurement Cycle}

Each measurement cycle consists of:
\begin{enumerate}
\item Lock cavity and lower atomic clock; record $R(h_1)$ for integration time $\tau$.
\item Reconfigure for upper station measurement.
\item Record $R(h_2)$ for integration time $\tau$.
\item Repeat with randomized order to decorrelate slow drifts.
\end{enumerate}

\paragraph{Integration budget.}
For $\tau \sim 10^4$ s per height and $N \sim 50$ cycles:
\begin{equation}
\sigma_{\mathrm{stat}} \sim \frac{10^{-18}}{\sqrt{N}} \sim 1.4 \times 10^{-19}.
\end{equation}
With systematic floor $\sigma_{\mathrm{syst}} \sim 2 \times 10^{-19}$:
\begin{equation}
\sigma_{\mathrm{tot}} = \sqrt{\sigma_{\mathrm{stat}}^2 + \sigma_{\mathrm{syst}}^2} 
\sim 2.5 \times 10^{-19}.
\end{equation}

\subsubsection{Systematics Budget}

\begin{table}[h]
\centering
\caption{Systematics budget for cavity-atom residual test.}\label{tab:cavity-systematics}
\scriptsize
\setlength{\tabcolsep}{2pt}
\begin{tabular}{lcc}
\toprule
Effect & Contrib. & Mitigation \\
\midrule
Temp.\ gradients & $< 10^{-19}$ & mK stab., shielding \\
Magnetic fields & $< 10^{-19}$ & nT stab., shielding \\
Pressure/refr.\ index & $< 10^{-20}$ & Vacuum \\
Vibrations & $< 10^{-20}$ & Isolation \\
Fiber link noise & $< 10^{-19}$ & Phase stab. \\
\midrule
\textbf{Total} & $\sim 2 \times 10^{-19}$ & \\
\bottomrule
\end{tabular}
\end{table}

\subsubsection{Blinding Protocol}

To prevent experimenter bias:
\begin{enumerate}
\item A secret offset $\delta$ at the $10^{-18}$ level is added to all 
recorded $R(h)$ values.
\item All data selection and systematic modeling performed on blinded data.
\item Analysis pipeline frozen before unblinding.
\item Offset removed only after all cuts finalized.
\end{enumerate}

\subsubsection{Pre-Registered Decision Rule}

Let $\widehat{\Delta R/R}$ be the unblinded estimator with uncertainty $\sigma_{\mathrm{tot}}$:

\begin{itemize}
\item \textbf{Null regime:} $|\widehat{\Delta R/R}| < 3\sigma_{\mathrm{tot}}$ 
$\Rightarrow$ consistent with GR and with geometric cancellation.  Upper bound on the residual:
\begin{equation}
|\xi_{\rm LPI}^{\rm res}| < \frac{3\sigma_{\mathrm{tot}}}{g\Delta h/c^2} \quad (95\%~\text{CL}).
\end{equation}

\item \textbf{Detection regime:} $|\widehat{\Delta R/R}| > 5\sigma_{\mathrm{tot}}$ 
$\Rightarrow$ a non-zero screened residual is measured:
\begin{equation}
\xi_{\rm LPI}^{\rm res} = \frac{\widehat{\Delta R/R}}{g\Delta h/c^2} \pm \frac{\sigma_{\mathrm{tot}}}{g\Delta h/c^2}.
\end{equation}

\item \textbf{Intermediate:} $3$--$5\sigma_{\mathrm{tot}}$ $\Rightarrow$ extend campaign.
\end{itemize}

\subsubsection{Sensitivity Reach}

For the benchmark parameters above, the minimum detectable residual is:
\begin{equation}
\xi_{\rm LPI,\,min}^{\rm res} \sim \frac{\sigma_{\mathrm{tot}}}{g\Delta h/c^2} \sim 2 \times 10^{-5}.
\end{equation}
This is sensitive enough to detect a residual at the level predicted by the screened formalism if it is near the upper end of the surviving window, but would require space-based or long-baseline platforms to push significantly deeper.

%-----------------------------------------------------------------------------
\subsection{Multi-Species Clock Comparison Protocol}\label{app:proto-clocks}
%-----------------------------------------------------------------------------

The full channel-resolved coupling of Eq.~(\ref{eq:K-general}) produces 
differential clock responses that can be measured without height separation.
The simplified pure-$\alpha$ scaling $K_A^{(\alpha)} = k_\alpha S_A^\alpha$
is only the leading same-ion term and is not the canonical master-law for
the v3.2 clock program.

\subsubsection{Observable}

For clock species $A$ and $B$ at the same location, measure:
\begin{equation}
\Delta_{AB}(t) \equiv \ln\frac{\nu_A(t)}{\nu_B(t)} - \langle \ln\frac{\nu_A}{\nu_B} \rangle.
\end{equation}

\paragraph{DFD prediction.}
Solar potential modulation at frequency $\Omega = 2\pi/\text{yr}$:
\begin{equation}
\Delta_{AB}(t) = (K_A - K_B) \cdot \frac{\Delta\Phi_\odot(t)}{c^2},
\end{equation}
where $\Delta\Phi_\odot/c^2 \sim 3 \times 10^{-10}$ over Earth's orbit.

\subsubsection{Species Selection}

\begin{table}[h]
\centering
\caption{Recommended clock species for DFD tests.}\label{tab:clock-species}
\begin{tabular}{llcc}
\toprule
Species & Transition & $S_A^\alpha$ & $K_A$ ($\times 10^{-5}$) \\
\midrule
Cs & Ground HFS & 2.83 & 2.83 \\
Rb & Ground HFS & 2.34 & 2.34 \\
H & 1S--2S & $\approx 0$ & $\approx 0$ \\
Sr & ${}^1S_0$--${}^3P_0$ & 0.06 & 0.06 \\
Yb & ${}^1S_0$--${}^3P_0$ & 0.31 & 0.31 \\
Al$^+$ & ${}^1S_0$--${}^3P_0$ & 0.008 & 0.008 \\
Hg$^+$ & ${}^2S_{1/2}$--${}^2D_{5/2}$ & $-3.2$ & $-3.2$ \\
Yb$^+$ (E3) & ${}^2S_{1/2}$--${}^2F_{7/2}$ & $-5.95$ & $-5.95$ \\
Th-229 & Nuclear & $\sim 10^4$ & $\sim 10$ \\
\bottomrule
\end{tabular}
\end{table}

\paragraph{Optimal pairs.}
Maximize $|S_A^\alpha - S_B^\alpha|$:
\begin{itemize}
\item Yb$^+$(E3)/H: $\Delta S \approx 6$
\item Yb$^+$(E3)/Al$^+$: $\Delta S \approx 6$
\item Cs/H: $\Delta S \approx 2.8$
\item Th-229/Sr: $\Delta S \sim 10^4$ (nuclear clock)
\end{itemize}

\subsubsection{Analysis Protocol}

\begin{enumerate}
\item Fit $\Delta_{AB}(t)$ to model: $A_0 + A_1\cos(\Omega t + \phi)$.
\item Extract amplitude $A_1$ and phase $\phi$.
\item Compare phase to predicted solar ephemeris.
\item If phase matches and $A_1 > 5\sigma$: detection.
\item If $A_1 < 3\sigma$: upper bound on $|K_A - K_B|$.
\end{enumerate}

%-----------------------------------------------------------------------------
\subsection{Matter-Wave Interferometry: $T^3$ Protocol}\label{app:proto-mwi}
%-----------------------------------------------------------------------------

Long-baseline matter-wave interferometers (MAGIS-100, AION) can detect the 
parity-isolated $T^3$ phase signature unique to DFD.

\subsubsection{Observable}

The DFD phase accumulation for interrogation time $T$:
\begin{equation}
\phi_{\mathrm{DFD}} = \phi_{\mathrm{GR}} + \delta\phi_{T^3},
\end{equation}
where the $T^3$ correction is:
\begin{equation}
\delta\phi_{T^3} = \eta_c \cdot k \cdot g \cdot T^3 \cdot \frac{a_*}{c^2},
\end{equation}
with $\eta_c = \alpha/4 \approx 1.8 \times 10^{-3}$.

\paragraph{Numerical estimate.}
For MAGIS-100 parameters ($T \sim 1$ s, $k \sim 10^7$ m$^{-1}$):
\begin{equation}
\delta\phi_{T^3} \sim 10^{-9}~\text{rad}.
\end{equation}

\subsubsection{Parity Isolation}

The $T^3$ term has \emph{opposite} parity under $g \to -g$ compared to the 
$T^2$ Newtonian term. This allows isolation via:
\begin{enumerate}
\item \textbf{Dual-launch:} Launch atoms up and down simultaneously.
\item \textbf{Differential measurement:} $\phi_{\mathrm{up}} - \phi_{\mathrm{down}}$.
\item \textbf{Result:} $T^2$ terms cancel; $T^3$ terms add.
\end{enumerate}

\subsubsection{Sensitivity Requirements}

\begin{table}[h]
\centering
\caption{MAGIS-100/AION sensitivity to DFD $T^3$ phase.}\label{tab:mwi-sensitivity}
\begin{tabular}{lccc}
\toprule
Facility & $T$ (s) & $\delta\phi_{T^3}$ (rad) & Detection threshold \\
\midrule
MAGIS-100 & 1.4 & $3 \times 10^{-9}$ & $10^{-10}$ rad/shot \\
AION-10 & 0.7 & $4 \times 10^{-10}$ & $10^{-11}$ rad/shot \\
AION-km & 2.3 & $2 \times 10^{-8}$ & $10^{-12}$ rad/shot \\
\bottomrule
\end{tabular}
\end{table}

\subsubsection{Falsification Criterion}

If parity-isolated $T^3$ phase is measured to be:
\begin{itemize}
\item Consistent with zero at $< 10^{-10}$ rad $\Rightarrow$ DFD $\eta_c$ prediction falsified.
\item Non-zero at $> 5\sigma$ $\Rightarrow$ New physics detected; DFD provides natural explanation.
\end{itemize}

%-----------------------------------------------------------------------------
\subsection{Nuclear Clock Protocol: Th-229}\label{app:proto-nuclear}
%-----------------------------------------------------------------------------

The ${}^{229}$Th nuclear isomer transition provides sensitivity to strong-sector 
couplings, with $d_s \sim 1.3$ (order of magnitude larger than $d_e$).

\subsubsection{Prediction}

DFD predicts:
\begin{equation}
K_{\mathrm{Th}} - K_{\mathrm{Sr}} \sim 8 \times 10^{-5},
\end{equation}
approximately $3\times$ larger than Cs/Sr difference.

\paragraph{Observable signal.}
For solar potential modulation:
\begin{equation}
\Delta\left[\ln\frac{\nu_{\mathrm{Th}}}{\nu_{\mathrm{Sr}}}\right] \sim 5 \times 10^{-15}.
\end{equation}

\subsubsection{Experimental Requirements}

\begin{itemize}
\item Nuclear clock operational with systematic uncertainty $< 10^{-16}$.
\item Continuous comparison with optical clock (Sr or Yb) over $\geq 1$ year.
\item Analysis for annual modulation at solar frequency.
\end{itemize}

\subsubsection{Timeline}

Nuclear clock technology is expected to reach required precision within 5--7 years.

%-----------------------------------------------------------------------------
\subsection{Space Mission Protocols}\label{app:proto-space}
%-----------------------------------------------------------------------------

Space-based tests provide access to larger potential differences and 
different systematic environments.

\subsubsection{ACES (ISS)}

The Atomic Clock Ensemble in Space provides:
\begin{itemize}
\item $\Delta\Phi/c^2 \sim 10^{-10}$ (ISS altitude).
\item Microwave clock comparisons with ground.
\item Sensitivity to $K_A - K_B$ at $10^{-7}$ level.
\end{itemize}

\subsubsection{Dedicated LPI Mission}

A dedicated mission with optical clocks could achieve:
\begin{itemize}
\item Highly elliptical orbit: $\Delta\Phi/c^2 \sim 10^{-9}$.
\item Cavity-atom comparison in space.
\item Sensitivity to screened residual $\xi_{\rm LPI}^{\rm res}$ at $10^{-6}$ level.
\end{itemize}

%-----------------------------------------------------------------------------
\subsection{Summary: Experimental Roadmap}\label{app:proto-summary}
%-----------------------------------------------------------------------------

\begin{table}[h]
\centering
\caption{DFD experimental verification timeline.}\label{tab:exp-roadmap}
\scriptsize
\setlength{\tabcolsep}{2pt}
\begin{tabular*}{\columnwidth}{@{\extracolsep{\fill}}llll@{}}
\toprule
Time & Test & Prediction & Falsification \\
\midrule
Now & UVCS & $\Gamma = 4$ & $\Gamma = 1$ at ${>}5\sigma$ \\
1--3\,yr & Cross-species clocks & Channel residuals & All-channel nulls \\
1--3\,yr & Nuclear clocks & 26\,Hz--kHz window & No annual signal \\
3--7\,yr & Matter-wave $T^3$ & $\delta\phi_{T^3} \neq 0$ & Null at $10^{-10}$ \\
${>}$7\,yr & Cavity--atom & Screened residual & Null at target \\
${>}$7\,yr & Space missions & Enhanced prec. & --- \\
\bottomrule
\end{tabular*}
\end{table}

\begin{tcolorbox}[colback=blue!5!white, colframe=blue!50!black, title=Experimental Protocol Summary]
\textbf{All protocols are pre-registered:}
\begin{itemize}[nosep]
\item Observables and predictions specified before data collection
\item Decision rules fixed in advance
\item Blinding protocols where applicable
\item Clear falsification criteria for both GR and DFD
\end{itemize}

\textbf{Key discriminators:}
\begin{itemize}[nosep]
\item Cavity--atom residual: screened non-metric mismatch after tree-level cancellation
\item Multi-species clocks: channel-resolved species dependence governed by Eq.~\eqref{eq:K-general}
\item Matter-wave $T^3$: Parity-isolated phase with DFD-specific scaling
\item Nuclear clocks: Strong-sector coupling $d_s \sim 1.3$
\end{itemize}

\textbf{Current status:}
\begin{itemize}[nosep]
\item UVCS double-transit: CONFIRMED ($\Gamma = 4.4 \pm 0.9$)
\item Others: Awaiting experimental implementation
\end{itemize}
\end{tcolorbox}
